\subsection{The two reciprocal-totient sums}
The prime and cofactor witnesses contribute equally to the leading
logarithmic mass. The following two estimates make that equality
explicit.

\begin{lemma}\label{lem:totient-three}
For $x\ge2$,
\[
 \sum_{\substack{m\le x\\m\equiv3\pmod4}}\frac1{\varphi(m)}
       =\frac{C_*}{2}\log x+O(1).
\]
\end{lemma}
\begin{proof}
Let $\chi$ be the nonprincipal character modulo four. The identity
\[
 \frac{n}{\varphi(n)}
   =\sum_{d\mid n}\frac{\mu^2(d)}{\varphi(d)}
\]
gives
\[
 H_{\rm odd}(x):=\sum_{\substack{n\le x\\n\ {\rm odd}}}
                      \frac1{\varphi(n)}
 =\sum_{\substack{d\le x\\d\ {\rm odd}}}
    \frac{\mu^2(d)}{d\varphi(d)}
    \sum_{\substack{k\le x/d\\k\ {\rm odd}}}\frac1k .
\]
The inner sum is $\frac12\log(x/d)+O(1)$ uniformly for $x/d\ge1$.
The coefficient sums
\[
 \sum_{d\ge1}\frac{\mu^2(d)}{d\varphi(d)}
 \quad\hbox{and}\quad
 \sum_{d\ge1}\frac{\mu^2(d)\log d}{d\varphi(d)}
\]
converge, using the elementary bound
$\varphi(d)\ge\sqrt{d/2}$.\footnote{For an odd prime power
$\ell^a$, $\varphi(\ell^a)/\sqrt{\ell^a}
=\ell^{a/2-1}(\ell-1)\ge1$.
For powers of two the ratio is at least $1/\sqrt2$.
Multiply these inequalities over the prime powers dividing $d$.}
The tail of the first sum multiplied by $\log x$
is bounded. It follows that
\[
 H_{\rm odd}(x)
 =\frac{\log x}{2}
   \prod_{\ell>2}\left(1+\frac1{\ell(\ell-1)}\right)+O(1)
 =C_*\log x+O(1).
\]
The product runs over primes and its coefficient equals $C_*$
because the missing prime-two factor in
$\prod_\ell(1+1/(\ell(\ell-1)))$ is $3/2$.

The same divisor identity, multiplied by $\chi(n)$, gives
\[
 \sum_{n\le x}\frac{\chi(n)}{\varphi(n)}
 =\sum_{d\le x}\frac{\mu^2(d)\chi(d)}{d\varphi(d)}
                \sum_{k\le x/d}\frac{\chi(k)}k=O(1).
\]
Indeed the character harmonic partial sums are uniformly bounded,
by summation by parts and the bounded partial sums of $\chi$,
and the outer coefficient sum converges absolutely.
For odd $m$, the indicator of $m\equiv3\pmod4$ is
$(1-\chi(m))/2$. Combining these two estimates proves the lemma.
\end{proof}

\begin{lemma}\label{lem:totient-even}
For $x\ge2$,
\[
 \sum_{\substack{i\le x\\i\ {\rm even}}}\frac1{\varphi(4i)}
       =\frac{C_*}{2}\log x+O(1).
\]
\end{lemma}
\begin{proof}
Write $i=2^a j$ with $a\ge1$ and $j$ odd. Since
$\varphi(4i)=2^{a+1}\varphi(j)$, the preceding proof gives
\[
 \sum_{\substack{i\le x\\i\ {\rm even}}}\frac1{\varphi(4i)}
 =\sum_{1\le a\le\log_2x}2^{-a-1}
       \bigl(C_*\log(x/2^a)+O(1)\bigr).
\]
The sums of $2^{-a-1}$ and $a2^{-a-1}$ over $a\ge1$
are $1/2$ and a finite constant respectively.
Extending these geometric sums to infinity changes the expression
by $O(1)$, proving the claim.
\end{proof}

\subsection{A uniform mass bound on fixed logarithmic intervals}
Set $w_f=\omega(f)/f$ for odd primes $f$.
\begin{lemma}\label{lem:short-block-mass}
For every $0<c<C_*/2$, all sufficiently small fixed $\delta>0$
have the following property. There is an integer $J\ge1$ such that
\begin{equation}\label{eq:log-block-mass}
 \sum_{\substack{f\ {\rm prime}\\
          e^{\delta j}<f\le e^{\delta(j+1)}}}w_f\ge c\delta
 \qquad(j\ge J).
\end{equation}
The parameters $c,\delta,J$ are fixed before the limit in $Q$.
\end{lemma}
\begin{proof}
Fix $0<\theta<1/2$ and put $D=\lfloor t^\theta\rfloor$.
For $t<f\le e^\delta t$, count the following divisors of $f+1$.

First count $m\le D$ with $m\equiv3\pmod4$ and
$f\equiv-1\pmod m$. Then count the divisors $m=(f+1)/i$
arising from even $i\le D$ and
\[
 f\equiv-1-i\pmod{4i}.
\]
The latter congruence gives $m\equiv3\pmod4$.
Moreover $m>t^{1-\theta}>D$, so the second collection is
disjoint from the first, and distinct $i$ give distinct divisors.

The residues $-1\pmod m$ and $-1-i\pmod{4i}$ are reduced,
the latter because $i$ is even.
Bombieri--Vinogradov, applied at both endpoints and with its
maximum over reduced residue classes, gives
\begin{align*}
 &\sum_{t<f\le e^\delta t}\omega(f)\\
 &\quad\ge
  \bigl(\li(e^\delta t)-\li(t)\bigr)
  \left(
     \sum_{\substack{m\le D\\m\equiv3\ (4)}}\frac1{\varphi(m)}
     +\sum_{\substack{i\le D\\i\ {\rm even}}}
                                    \frac1{\varphi(4i)}
  \right)
  +O_{\theta,\delta}\bigl(t/(\log t)^2\bigr).
\end{align*}
Every modulus is at most $4t^\theta$, which lies below the usual
level $t^{1/2}(\log t)^{-A}$ for sufficiently large $t$.
The sign of the error above is immaterial to this lower bound.

Lemmas~\ref{lem:totient-three} and~\ref{lem:totient-even}
make the parenthesis $C_*\theta\log t+O(1)$.
For fixed $\delta$, the logarithmic integral difference is
$(e^\delta-1)t/\log t+O_\delta(t/(\log t)^2)$.
Since $1/f\ge e^{-\delta}/t$ on the interval, we obtain
\begin{equation}\label{eq:short-block-prelimit}
 \sum_{t<f\le e^\delta t}w_f
 \ge C_*\theta(1-e^{-\delta})-o_{\theta,\delta}(1).
\end{equation}
Choose $\theta$ so that $c<C_*\theta<C_*/2$.
As $\delta$ tends to zero,
$C_*\theta(1-e^{-\delta})/\delta$ tends to $C_*\theta$.
Thus every sufficiently small fixed $\delta$ makes that quotient
strictly greater than $c$. For that choice, take $t=e^{\delta j}$
and then $J$ sufficiently large to absorb the error uniformly for
all $j\ge J$. This proves \eqref{eq:log-block-mass}.
\end{proof}

\subsection{Products from distinct logarithmic blocks}
The use of distinct prime blocks gives a direct lower bound for
the large-sieve denominator without repeated prime factors.

\begin{proposition}\label{prop:sharper-L}
One has
\begin{equation}\label{eq:sharper-L}
 \liminf_{Q\to\infty}\frac{\log L(Q)}{\sqrt{\log Q}}
       \ge\sqrt{2C_*}.
\end{equation}
Equivalently,
$\log L(Q)\ge(\sqrt{2C_*}-o(1))\sqrt{\log Q}$.
\end{proposition}
\begin{proof}
Fix $0<c<C_*/2$ and $0<\eps<1$. Choose a fixed $\delta>0$
small enough for Lemma~\ref{lem:short-block-mass} and for
$c\delta<\eps$. Let $J$ be its threshold.
Put
\[
 U=\log Q,\qquad
 k=\left\lfloor\sqrt{c(1-\eps)U}\right\rfloor,\qquad
 R=\lfloor U/\delta\rfloor-k(J+1)-\frac{k(k-1)}2.
\]
For sufficiently large $Q$, the integer $R$ is positive.
Choose $k$ distinct block indices
$J\le j_1<\cdots<j_k$ and one prime from each chosen block.
Their product is at most $Q$ whenever
$\delta\sum_{r=1}^k(j_r+1)\le U$.
The primes are distinct, and unique factorization assigns every
resulting squarefree integer to just one set of prime factors
and just one set of blocks.

Write
\[
 j_r=J+(r-1)+b_r,\qquad 0\le b_1\le\cdots\le b_k .
\]
It suffices that $\sum b_r\le R$. There are
$\binom{R+k}{k}$ ordered nonnegative integer vectors with
coordinate sum at most $R$. Sorting sends them to nondecreasing
vectors with at most $k!$ preimages for each sorted vector.
Therefore at least $\binom{R+k}{k}/k!$ block selections satisfy
the product condition. In $L(Q)$ the prime factor weight is
$\omega(f)/(f-\omega(f))\ge w_f$.
Summing independently over one prime in each selected block
and applying \eqref{eq:log-block-mass} gives
\begin{equation}\label{eq:distinct-block-bound}
 L(Q)\ge (c\delta)^k\frac{\binom{R+k}{k}}{k!}
       \ge\frac{(c\delta R)^k}{(k!)^2}.
\end{equation}
Here the second inequality follows by expanding the binomial
coefficient as $\prod_{r=1}^k(R+r)/r$.

Our choices imply
\[
 \delta R
 \ge U-\delta-\delta k(J+1)-\frac{\delta k(k-1)}2
 \ge (1-\eps)U
\]
for all sufficiently large $Q$.
Indeed the triangular term is at most $\eps U/2$, and the
other losses are $O_{c,\delta,J,\eps}(\sqrt U)$.
Set $B=c(1-\eps)U$. Then $k=\lfloor\sqrt B\rfloor$ and
\eqref{eq:distinct-block-bound}, followed by Stirling's formula,
gives
\[
 \log L(Q)\ge k\log B-2\log(k!)
 \ge 2\sqrt{c(1-\eps)U}-O_{c,\eps}(\log U).
\]
Taking a limit inferior first, and then letting $c$ increase
to $C_*/2$ and $\eps$ decrease to zero, proves the result.
\end{proof}

\subsection{The saving for square-root search depth}
We now complete the growing-depth conclusion of
Theorem~\ref{thm:large} with the strengthened constant.
Discard primes $p\le\sqrt X$ and use dyadic intervals $(Y/2,Y]$
for the rest. Put
\[
 M_Y=\lfloor\sqrt{Y/2}/2\rfloor,\qquad Q=3M_Y-1.
\]
Then $M(p)\ge M_Y$ in the interval,
$Q^2\le9Y/8$, and $\log Q=\frac12\log Y+O(1)$.
For any fixed $0<a<\sqrt{C_*}$,
Proposition~\ref{prop:sharper-L} and the arithmetic
large-sieve inequality therefore give
\[
 \#\{Y/2<p\le Y:p\in\bigcap_{i\le M(p)}F_i\}
       \ll_a Y\exp(-a\sqrt{\log Y}).
\]
For $Y=X/2^j$, use
\[
 \sqrt{\log X-j\log2}
 \ge\sqrt{\log X}-\frac{j\log2}{\sqrt{\log X}}
\]
whenever $0\le j\log2\le\log X$.
The dyadic sum is bounded by $Xe^{-a\sqrt{\log X}}$
times a geometric series of bounded sum for sufficiently large
$X$. The discarded $O(\sqrt X)$ primes are absorbed. We obtain
\begin{equation}\label{eq:sharper-sqrt}
 E_{\sqrt{\ }}(X)\ll_a X\exp(-a\sqrt{\log X})
 \qquad(0<a<\sqrt{C_*}).
\end{equation}
The threshold furnished here is
\[
 \sqrt{C_*}=0.804900912498\ldots,\qquad
 \sqrt{2C_*}=1.138301786821\ldots .
\]
It concerns the failure of the stated source budget. The following
pointwise analysis records an explicit failure of that budget
and the exact finite search needed for the full equation.
